\subsection{PROOFS}

A demonstrative text in a theory \(\mathfrak{G}\) comprises:

\begin{enumerate}
\def\labelenumi{(\arabic{enumi})}
\item
  An auxiliary formative construction of relations and terms in \(\mathfrak{G}\) .
\item
  A proof in \(\mathfrak{G}\) , that is to say a sequence of relations in \(\mathfrak{G}\) which appear in the auxiliary formative construction, such that for every relation \(\pmb{R}\) in the sequence at least one of the following conditions is satisfied:
\end{enumerate}

\((a_{1})\boldsymbol{R}\) is an explicit axiom of \(\mathfrak{c}\) .

\((a_{2})\) R results from the application of a scheme of C to terms or relations which appear in the auxiliary formative construction.

\begin{enumerate}
\def\labelenumi{(\alph{enumi})}
\setcounter{enumi}{1}
\tightlist
\item
  there are two relations \(S, T\) in the sequence which precede \(R\) , such that \(T\) is \(S \Longrightarrow R\) .
\end{enumerate}

¶ A theorem in \(\mathfrak{G}\) is a relation which appears in a proof in \(\mathfrak{G}\).

This notion is therefore essentially dependent on the state of the theory under consideration, at the time when it is being described. A relation in a theory \(\mathfrak{G}\) becomes a theorem in \(\mathfrak{G}\) when one succeeds in inserting it into a proof in \(\mathfrak{G}\) . To say that a relation in \(\mathfrak{G}\) ``is not a theorem in \(\mathfrak{G}\)'' cannot have any meaning without reference to the stage of development of the theory \(\mathfrak{G}\) .

A theorem in \(\mathfrak{G}\) is also called a ``true relation in \(\mathfrak{G}\) '' (or ``proposition'', ``lemma'', ``corollary'', etc.). Let \(R\) be a relation in \(\mathfrak{G}\) , let \(x\) be a letter and \(T\) a term in \(\mathfrak{G}\) ; if \((T|x)R\) is a theorem in \(\mathfrak{G}\) , \(T\) is said to satisfy the relation \(R\) in \(\mathfrak{G}\) (or to be a solution of \(R\) ), when \(R\) is considered as a relation in \(x\) .

A relation is said to be false in \(\mathfrak{c}\) if its negation is a theorem in \(\mathfrak{c}\) . A theory \(\mathfrak{c}\) is said to be contradictory when one has written a relation which is both true and false in \(\mathfrak{c}\) .

Here again, we are dealing with a notion that depends on the particular state of development of a theory. The reader should beware of the confusion (unfortunately suggested by the intuitive meaning of the word ``false'') which consists in believing that, once one has proved that a relation R is false in \(\mathfrak{G}\) , one has thereby established that R ``is not true'' in \(\mathfrak{G}\) (strictly speaking, this phrase has no precise meaning in mathematics, as we have remarked above).

¶ In what follows we shall give metamathematical criteria, called deductive criteria, which will allow us to shorten proofs. These criteria will be denoted by the letter C followed by a number.

C1 (Syllogism). Let A and B be relations in a theory C. If A and \(A \Longrightarrow B\) are theorems in C, then B is a theorem in C.

Let \(R_1, R_2, \ldots, R_n\) be a proof in \(\mathfrak{G}\) in which \(A\) appears, and let \(S_1, S_2, \ldots, S_p\) be a proof in \(\mathfrak{G}\) in which \(A \Longrightarrow B\) appears. Clearly \(R_1, R_2, \ldots, R_n, S_1, S_2, \ldots, S_p\) is a proof in \(\mathfrak{G}\) in which both \(A\) and \(A \Longrightarrow B\) appear. Hence

\[
\boldsymbol {R} _ {1}, \boldsymbol {R} _ {2}, \dots , \boldsymbol {R} _ {n}, \boldsymbol {S} _ {1}, \boldsymbol {S} _ {2}, \dots , \boldsymbol {S} _ {p}, \boldsymbol {B}
\]

is a proof in \(\mathfrak{c}\) , and therefore \(\pmb{B}\) is a theorem in \(\mathfrak{c}\) .

